\subsubsection{More about lengths and simples}

Let us continue studying lengths of permutations.

\begin{proposition}
\label{prop.perm.len.inv}Let $n\in\mathbb{N}$ and $\sigma\in S_{n}$. Then,
$\ell\left(  \sigma^{-1}\right)  =\ell\left(  \sigma\right)  $.
\end{proposition}

\begin{proof}
[Proof of Proposition \ref{prop.perm.len.inv} (sketched).](See \cite[Exercise
5.2 \textbf{(f)}]{detnotes} for details.)

Recall that $\ell\left(  \sigma\right)  $ is the \# of inversions of $\sigma$,
while $\ell\left(  \sigma^{-1}\right)  $ is the \# of inversions of
$\sigma^{-1}$.

Recall also that an inversion of $\sigma$ is a pair $\left(  i,j\right)
\in\left[  n\right]  \times\left[  n\right]  $ such that $i<j$ and
$\sigma\left(  i\right)  >\sigma\left(  j\right)  $. Likewise, an inversion of
$\sigma^{-1}$ is a pair $\left(  u,v\right)  \in\left[  n\right]
\times\left[  n\right]  $ such that $u<v$ and $\sigma^{-1}\left(  u\right)
>\sigma^{-1}\left(  v\right)  $.

Thus, if $\left(  i,j\right)  $ is an inversion of $\sigma$, then $\left(
\sigma\left(  j\right)  ,\sigma\left(  i\right)  \right)  $ is an inversion of
$\sigma^{-1}$. Hence, we obtain a map%
\begin{align*}
\left\{  \text{inversions of }\sigma\right\}   &  \rightarrow\left\{
\text{inversions of }\sigma^{-1}\right\}  ,\\
\left(  i,j\right)   &  \mapsto\left(  \sigma\left(  j\right)  ,\sigma\left(
i\right)  \right)  .
\end{align*}
This map is furthermore bijective (indeed, it has an inverse map, which sends
each $\left(  u,v\right)  \in\left\{  \text{inversions of }\sigma
^{-1}\right\}  $ to $\left(  \sigma^{-1}\left(  v\right)  ,\sigma^{-1}\left(
u\right)  \right)  $). Thus, the bijection principle yields%
\[
\left(  \text{\# of inversions of }\sigma\right)  =\left(  \text{\# of
inversions of }\sigma^{-1}\right)  .
\]
In other words, $\ell\left(  \sigma\right)  =\ell\left(  \sigma^{-1}\right)
$. This proves Proposition \ref{prop.perm.len.inv}.
\end{proof}

The following lemma is crucial for understanding lengths of permutations:

\begin{lemma}
[single swap lemma]\label{lem.perm.len.ssl}Let $n\in\mathbb{N}$, $\sigma\in
S_{n}$ and $k\in\left[  n-1\right]  $. Then: \medskip

\textbf{(a)} We have%
\[
\ell\left(  \sigma s_{k}\right)  =%
\begin{cases}
\ell\left(  \sigma\right)  +1, & \text{if }\sigma\left(  k\right)
<\sigma\left(  k+1\right)  ;\\
\ell\left(  \sigma\right)  -1, & \text{if }\sigma\left(  k\right)
>\sigma\left(  k+1\right)  .
\end{cases}
\]


\textbf{(b)} We have%
\[
\ell\left(  s_{k}\sigma\right)  =%
\begin{cases}
\ell\left(  \sigma\right)  +1, & \text{if }\sigma^{-1}\left(  k\right)
<\sigma^{-1}\left(  k+1\right)  ;\\
\ell\left(  \sigma\right)  -1, & \text{if }\sigma^{-1}\left(  k\right)
>\sigma^{-1}\left(  k+1\right)  .
\end{cases}
\]


[\textbf{Note:} If $i\in\left[  n\right]  $, then $\sigma\left(  i\right)  $
is the \textbf{entry} in position $i$ of the one-line notation of $\sigma$,
whereas $\sigma^{-1}\left(  i\right)  $ is the \textbf{position} in which the
number $i$ appears in the one-line notation of $\sigma$. For example, if
$\sigma=512364$ in one-line notation, then $\sigma\left(  6\right)  =4$ and
$\sigma^{-1}\left(  6\right)  =5$.]
\end{lemma}

We will only outline the proof of Lemma \ref{lem.perm.len.ssl}; a detailed
proof can be found in \cite[Exercise 5.2 \textbf{(a)}]{detnotes} (although it
is a slightly different proof).

\begin{proof}
[Proof of Lemma \ref{lem.perm.len.ssl} (sketched).]\textbf{(b)} The
OLN\footnote{Recall that \textquotedblleft OLN\textquotedblright\ means
\textquotedblleft one-line notation\textquotedblright.} of $s_{k}\sigma$ is
obtained from the OLN of $\sigma$ by swapping the two entries $k$ and $k+1$.
This is best seen on an example: For example, if $\sigma=512364$ (in OLN),
then $s_{3}\sigma=512463$. In general, this follows by observing that%
\[
\left(  s_{k}\sigma\right)  \left(  i\right)  =s_{k}\left(  \sigma\left(
i\right)  \right)  =%
\begin{cases}
k+1, & \text{if }\sigma\left(  i\right)  =k;\\
k, & \text{if }\sigma\left(  i\right)  =k+1;\\
\sigma\left(  i\right)  , & \text{otherwise}%
\end{cases}
\ \ \ \ \ \ \ \ \ \ \text{for each }i\in\left[  n\right]  .
\]


Let us now use this observation to see how the inversions of $s_{k}\sigma$
differ from the inversions of $\sigma$. Indeed, let us call an inversion
$\left(  i,j\right)  $ of a permutation $\tau$ \emph{exceptional} if we have
$\tau\left(  i\right)  =k+1$ and $\tau\left(  j\right)  =k$. All other
inversions of $\tau$ will be called \emph{non-exceptional}.

Now, we make the following observation:

\begin{statement}
\textit{Observation 1:} If $\left(  i,j\right)  $ is any non-exceptional
inversion of $\sigma$, then $\left(  i,j\right)  $ is still a non-exceptional
inversion of $s_{k}\sigma$.
\end{statement}

[\textit{Proof of Observation 1:} Let $\left(  i,j\right)  $ be a
non-exceptional inversion of $\sigma$. Thus, we have the inequality
$\sigma\left(  i\right)  >\sigma\left(  j\right)  $ (since $\left(
i,j\right)  $ is an inversion of $\sigma$). This inequality cannot get
reversed by applying $s_{k}$ to both its sides (i.e., we cannot have
$s_{k}\left(  \sigma\left(  i\right)  \right)  <s_{k}\left(  \sigma\left(
j\right)  \right)  $), since the only pair $\left(  u,v\right)  \in\left[
n\right]  \times\left[  n\right]  $ satisfying $u>v$ and $s_{k}\left(
u\right)  <s_{k}\left(  v\right)  $ is the pair $\left(  k+1,k\right)  $ (but
our pair $\left(  \sigma\left(  i\right)  ,\sigma\left(  j\right)  \right)  $
cannot equal this pair $\left(  k+1,k\right)  $, since the inversion $\left(
i,j\right)  $ of $\sigma$ is non-exceptional). Hence, we have $s_{k}\left(
\sigma\left(  i\right)  \right)  \geq s_{k}\left(  \sigma\left(  j\right)
\right)  $. In other words, $\left(  s_{k}\sigma\right)  \left(  i\right)
\geq\left(  s_{k}\sigma\right)  \left(  j\right)  $. Since the map
$s_{k}\sigma$ is injective (being a permutation of $\left[  n\right]  $), we
thus obtain $\left(  s_{k}\sigma\right)  \left(  i\right)  >\left(
s_{k}\sigma\right)  \left(  j\right)  $ (since $i\neq j$). Thus, the pair
$\left(  i,j\right)  $ is an inversion of $s_{k}\sigma$. Moreover, this
inversion $\left(  i,j\right)  $ is non-exceptional (since otherwise we would
have $\left(  s_{k}\sigma\right)  \left(  i\right)  =k+1$ and $\left(
s_{k}\sigma\right)  \left(  j\right)  =k$, which would lead to $\sigma\left(
i\right)  =k$ and $\sigma\left(  j\right)  =k+1$, which would contradict
$\sigma\left(  i\right)  >\sigma\left(  j\right)  $). Thus, we have shown that
$\left(  i,j\right)  $ is still a non-exceptional inversion of $s_{k}\sigma$.
This proves Observation 1.]

Similarly to Observation 1, we can prove the following:

\begin{statement}
\textit{Observation 2:} If $\left(  i,j\right)  $ is any non-exceptional
inversion of $s_{k}\sigma$, then $\left(  i,j\right)  $ is still a
non-exceptional inversion of $\sigma$.
\end{statement}

(Alternatively, we can obtain Observation 2 by applying Observation 1 to
$s_{k}\sigma$ instead of $\sigma$, since we have $\underbrace{s_{k}s_{k}%
}_{=s_{k}^{2}=\operatorname*{id}}\sigma=\sigma$.)

Combining Observation 1 with Observation 2, we see that the non-exceptional
inversions of $s_{k}\sigma$ are precisely the non-exceptional inversions of
$\sigma$. Hence,%
\begin{align}
&  \left(  \text{\# of non-exceptional inversions of }s_{k}\sigma\right)
\nonumber\\
&  =\left(  \text{\# of non-exceptional inversions of }\sigma\right)  .
\label{pf.lem.perm.len.ssl.eq-nonexc}%
\end{align}


What about the exceptional inversions? A permutation $\tau\in S_{n}$ has a
unique exceptional inversion if $k$ appears after $k+1$ in the OLN of $\tau$
(that is, if we have $\tau^{-1}\left(  k\right)  >\tau^{-1}\left(  k+1\right)
$); otherwise, it has none. Thus:

\begin{itemize}
\item If $\sigma^{-1}\left(  k\right)  <\sigma^{-1}\left(  k+1\right)  $, then
the permutation $s_{k}\sigma$ has a unique exceptional inversion, whereas the
permutation $\sigma$ has none.

\item If $\sigma^{-1}\left(  k\right)  >\sigma^{-1}\left(  k+1\right)  $, then
the permutation $\sigma$ has a unique exceptional inversion, whereas the
permutation $s_{k}\sigma$ has none.
\end{itemize}

Thus,%
\begin{align}
&  \left(  \text{\# of exceptional inversions of }s_{k}\sigma\right)
\nonumber\\
&  =\left(  \text{\# of exceptional inversions of }\sigma\right) \nonumber\\
&  \ \ \ \ \ \ \ \ \ \ +%
\begin{cases}
1, & \text{if }\sigma^{-1}\left(  k\right)  <\sigma^{-1}\left(  k+1\right)
;\\
-1, & \text{if }\sigma^{-1}\left(  k\right)  >\sigma^{-1}\left(  k+1\right)  .
\end{cases}
\label{pf.lem.perm.len.ssl.eq-exc}%
\end{align}


Now, recall that each inversion of a permutation $\tau\in S_{n}$ is either
exceptional or non-exceptional (and cannot be both at the same time). Thus,
adding together the two equalities (\ref{pf.lem.perm.len.ssl.eq-exc}) and
(\ref{pf.lem.perm.len.ssl.eq-nonexc}), we obtain%
\begin{align*}
&  \left(  \text{\# of inversions of }s_{k}\sigma\right) \\
&  =\left(  \text{\# of inversions of }\sigma\right)  +%
\begin{cases}
1, & \text{if }\sigma^{-1}\left(  k\right)  <\sigma^{-1}\left(  k+1\right)
;\\
-1, & \text{if }\sigma^{-1}\left(  k\right)  >\sigma^{-1}\left(  k+1\right)
\end{cases}
\\
&  =%
\begin{cases}
\left(  \text{\# of inversions of }\sigma\right)  +1, & \text{if }\sigma
^{-1}\left(  k\right)  <\sigma^{-1}\left(  k+1\right)  ;\\
\left(  \text{\# of inversions of }\sigma\right)  -1, & \text{if }\sigma
^{-1}\left(  k\right)  >\sigma^{-1}\left(  k+1\right)  .
\end{cases}
\end{align*}
In other words,%
\[
\ell\left(  s_{k}\sigma\right)  =%
\begin{cases}
\ell\left(  \sigma\right)  +1, & \text{if }\sigma^{-1}\left(  k\right)
<\sigma^{-1}\left(  k+1\right)  ;\\
\ell\left(  \sigma\right)  -1, & \text{if }\sigma^{-1}\left(  k\right)
>\sigma^{-1}\left(  k+1\right)
\end{cases}
\]
(since $\ell\left(  \tau\right)  $ denotes the \# of inversions of any
permutation $\tau$). This proves Lemma \ref{lem.perm.len.ssl} \textbf{(b)}.
\medskip

\textbf{(a)} Applying Lemma \ref{lem.perm.len.ssl} \textbf{(b)} to
$\sigma^{-1}$ instead of $\sigma$, we obtain%
\begin{align}
\ell\left(  s_{k}\sigma^{-1}\right)   &  =%
\begin{cases}
\ell\left(  \sigma^{-1}\right)  +1, & \text{if }\left(  \sigma^{-1}\right)
^{-1}\left(  k\right)  <\left(  \sigma^{-1}\right)  ^{-1}\left(  k+1\right)
;\\
\ell\left(  \sigma^{-1}\right)  -1, & \text{if }\left(  \sigma^{-1}\right)
^{-1}\left(  k\right)  >\left(  \sigma^{-1}\right)  ^{-1}\left(  k+1\right)
\end{cases}
\nonumber\\
&  =%
\begin{cases}
\ell\left(  \sigma\right)  +1, & \text{if }\sigma\left(  k\right)
<\sigma\left(  k+1\right)  ;\\
\ell\left(  \sigma\right)  -1, & \text{if }\sigma\left(  k\right)
>\sigma\left(  k+1\right)
\end{cases}
\label{pf.lem.perm.len.ssl.a.1}%
\end{align}
(since $\left(  \sigma^{-1}\right)  ^{-1}=\sigma$, and since Proposition
\ref{prop.perm.len.inv} yields $\ell\left(  \sigma^{-1}\right)  =\ell\left(
\sigma\right)  $). However, Proposition \ref{prop.perm.len.inv} (applied to
$\sigma s_{k}$ instead of $\sigma$) yields $\ell\left(  \left(  \sigma
s_{k}\right)  ^{-1}\right)  =\ell\left(  \sigma s_{k}\right)  $. In view of
$\left(  \sigma s_{k}\right)  ^{-1}=\underbrace{s_{k}^{-1}}_{=s_{k}}%
\sigma^{-1}=s_{k}\sigma^{-1}$, this rewrites as $\ell\left(  s_{k}\sigma
^{-1}\right)  =\ell\left(  \sigma s_{k}\right)  $. Comparing this with
(\ref{pf.lem.perm.len.ssl.a.1}), we obtain%
\[
\ell\left(  \sigma s_{k}\right)  =%
\begin{cases}
\ell\left(  \sigma\right)  +1, & \text{if }\sigma\left(  k\right)
<\sigma\left(  k+1\right)  ;\\
\ell\left(  \sigma\right)  -1, & \text{if }\sigma\left(  k\right)
>\sigma\left(  k+1\right)  .
\end{cases}
\]
This proves Lemma \ref{lem.perm.len.ssl} \textbf{(a)}.
\end{proof}

Lemma \ref{lem.perm.len.ssl} answers what happens to the length of a
permutation when we compose it (from the left or the right) with a simple
transposition $s_{k}$. What happens when we compose it with a non-simple
transposition? The situation is more complicated, but it is still true that
the length decreases or increases depending on whether the two entries that
are being swapped formed an inversion or not. Here is the exact answer (stated
only for $\sigma t_{i,j}$, but a version for $t_{i,j}\sigma$ can easily be
derived from it):

\begin{proposition}
\label{prop.perm.lisitij}Let $n\in\mathbb{N}$ and $\sigma\in S_{n}$. Let $i$
and $j$ be two elements of $\left[  n\right]  $ such that $i<j$. Then:
\medskip

\textbf{(a)} We have $\ell\left(  \sigma t_{i,j}\right)  <\ell\left(
\sigma\right)  $ if $\sigma\left(  i\right)  >\sigma\left(  j\right)  $. We
have $\ell\left(  \sigma t_{i,j}\right)  >\ell\left(  \sigma\right)  $ if
$\sigma\left(  i\right)  <\sigma\left(  j\right)  $. \medskip

\textbf{(b)} We have%
\[
\ell\left(  \sigma t_{i,j}\right)  =%
\begin{cases}
\ell\left(  \sigma\right)  -2\left\vert Q\right\vert -1, & \text{if }%
\sigma\left(  i\right)  >\sigma\left(  j\right)  ;\\
\ell\left(  \sigma\right)  +2\left\vert R\right\vert +1, & \text{if }%
\sigma\left(  i\right)  <\sigma\left(  j\right)  ,
\end{cases}
\]
where
\begin{align*}
Q  &  =\left\{  k\in\left\{  i+1,i+2,\ldots,j-1\right\}  \ \mid\ \sigma\left(
i\right)  >\sigma\left(  k\right)  >\sigma\left(  j\right)  \right\}
\ \ \ \ \ \ \ \ \ \ \text{and}\\
R  &  =\left\{  k\in\left\{  i+1,i+2,\ldots,j-1\right\}  \ \mid\ \sigma\left(
i\right)  <\sigma\left(  k\right)  <\sigma\left(  j\right)  \right\}  .
\end{align*}

\end{proposition}

\begin{proof}
[Proof of Proposition \ref{prop.perm.lisitij} (sketched).]\textbf{(b)} This
follows by a diligent analysis of the possible interactions between an
inversion and composition by $t_{i,j}$. To be more concrete:

\begin{itemize}
\item The fact that $\ell\left(  \sigma t_{i,j}\right)  =\ell\left(
\sigma\right)  -2\left\vert Q\right\vert -1$ when $\sigma\left(  i\right)
>\sigma\left(  j\right)  $ is \cite[Exercise 5.20]{detnotes}. A
straightforward proof was given by Elafandi in \cite[solution to Exercise
1]{18f-hw4se}. (The proof given in \cite[solution to Exercise 5.20]{detnotes}
is more circuitous, as it uses summation tricks to bypass case distinctions.)

\item The fact that $\ell\left(  \sigma t_{i,j}\right)  =\ell\left(
\sigma\right)  +2\left\vert R\right\vert +1$ when $\sigma\left(  i\right)
<\sigma\left(  j\right)  $ follows by applying the previous fact to $\sigma
t_{i,j}$ instead of $\sigma$. (Indeed, if $\sigma\left(  i\right)
<\sigma\left(  j\right)  $, then $\left(  \sigma t_{i,j}\right)  \left(
i\right)  =\sigma\left(  j\right)  >\sigma\left(  i\right)  =\left(  \sigma
t_{i,j}\right)  \left(  j\right)  $ and $\sigma\underbrace{t_{i,j}t_{i,j}%
}_{=t_{i,j}^{2}=\operatorname*{id}}=\sigma$. Moreover, when we replace
$\sigma$ by $\sigma t_{i,j}$, the sets $Q$ and $R$ trade places.)
\end{itemize}

\textbf{(a)} This follows immediately from part \textbf{(b)}.
\end{proof}

Now we come to one of the main facts about permutations of a finite set.

\begin{convention}
\label{conv.perm.simple}We recall that a \emph{simple transposition} in
$S_{n}$ means one of the $n-1$ transpositions $s_{1},s_{2},\ldots,s_{n-1}$. We
shall occasionally abbreviate \textquotedblleft simple
transposition\textquotedblright\ as \textquotedblleft\emph{simple}%
\textquotedblright.
\end{convention}

\begin{theorem}
[1st reduced word theorem for the symmetric group]%
\label{thm.perm.len.redword1}Let $n\in\mathbb{N}$ and $\sigma\in S_{n}$. Then:
\medskip

\textbf{(a)} We can write $\sigma$ as a composition (i.e., product) of
$\ell\left(  \sigma\right)  $ simples. \medskip

\textbf{(b)} The number $\ell\left(  \sigma\right)  $ is the smallest
$p\in\mathbb{N}$ such that we can write $\sigma$ as a composition of $p$
simples. \medskip

[\textbf{Keep in mind:} The composition of $0$ simples is $\operatorname*{id}%
$, since $\operatorname*{id}$ is the neutral element of the group $S_{n}$.]
\end{theorem}

\begin{example}
Let $\sigma\in S_{4}$ be the permutation $4132$ (in OLN). How can we write
$\sigma$ as a composition of simples? There are several ways to do this; for
example,%
\[
\sigma=\underbrace{s_{2}s_{3}s_{2}}_{=s_{3}s_{2}s_{3}}s_{1}=s_{3}%
s_{2}\underbrace{s_{3}s_{1}}_{=s_{1}s_{3}}=s_{3}s_{2}s_{1}s_{3}=s_{2}%
s_{1}s_{1}s_{3}s_{2}s_{1}=s_{2}s_{1}s_{3}s_{1}s_{2}s_{1}=\cdots.
\]
The shortest of these representations involve $4$ simples, precisely as
predicted by Theorem \ref{thm.perm.len.redword1} (since $\ell\left(
\sigma\right)  =4$).
\end{example}

Before we prove Theorem \ref{thm.perm.len.redword1}, let me mention a
geometric visualization of the symmetric group that will not be used in what
follows, but sheds some light on the theorem and on the role of simple transpositions:

\begin{remark}
Let $n\in\mathbb{N}$. Then, the permutations of $\left[  n\right]  $ can be
represented as the vertices of an $\left(  n-1\right)  $-dimensional polytope
in $n$-dimensional space.

Namely, each permutation $\sigma$ of $\left[  n\right]  $ gives rise to the
point%
\[
V_{\sigma}:=\left(  \sigma\left(  1\right)  ,\sigma\left(  2\right)
,\ldots,\sigma\left(  n\right)  \right)  \in\mathbb{R}^{n}.
\]
The convex hull of all these $n!$ many points $V_{\sigma}$ (for $\sigma\in
S_{n}$) is a polytope (i.e., a bounded convex polyhedron in $\mathbb{R}^{n}$).
This polytope is known as the
\emph{\href{https://en.wikipedia.org/wiki/Permutohedron}{\emph{permutahedron}%
}} (corresponding to $n$). It is actually $\left(  n-1\right)  $-dimensional,
since all its vertices lie on the hyperplane with equation $x_{1}+x_{2}%
+\cdots+x_{n}=1+2+\cdots+n$. It can be shown (see, e.g., \cite{GaiGup77}) that:

\begin{itemize}
\item The vertices of this polytope are precisely the $n!$ points $V_{\sigma}$
with $\sigma\in S_{n}$.

\item Two vertices $V_{\sigma}$ and $V_{\tau}$ are joined by an edge if and
only if $\sigma=s_{k}\tau$ for some $k\in\left[  n-1\right]  $.
\end{itemize}

\noindent The (intuitively obvious) fact that any two vertices of a polytope
can be connected by a sequence of edges therefore yields that any $\sigma\in
S_{n}$ can be written as a product of simples. This is a weaker version of
Theorem \ref{thm.perm.len.redword1} \textbf{(a)}.

We refer to textbooks on discrete geometry and geometric combinatorics for
more about polytopes and permutahedra in particular. Let me here just show the
permutahedra for $n=3$ and for $n=4$ (note that the permutahedron for $n=2$ is
a boring line segment in $\mathbb{R}^{2}$):

\begin{itemize}
\item The permutahedron for $n=3$ is a regular hexagon:%
\[%
\begin{tabular}
[c]{ccc}%
$%
% [tikz figure removed]
%BeginExpansion
% [tikz figure removed]%
%EndExpansion
$ & $\ \ \ \ \ \ \ \ \ \ $ & $%
% [tikz figure removed]
%BeginExpansion
% [tikz figure removed]%
%EndExpansion
$%
\end{tabular}
\
\]
The picture on the left
(\href{https://tex.stackexchange.com/a/321238/}{courtesy of tex.stackexchange
user Jake}, released under the MIT license) shows the permutahedron embedded
in $\mathbb{R}^{3}$; the picture on the right is a view from an orthogonal direction.

\item The permutahedron for $n=4$ is a truncated octahedron:%
\[%
% [svg figure removed]
%BeginExpansion
\incpdftexpic{Permutohedron.pdf_tex}%
%EndExpansion
\]
(\href{https://en.wikipedia.org/wiki/File:Permutohedron.svg}{picture courtesy
of David Eppstein on the Wikipedia}).
\end{itemize}
\end{remark}

Let us now outline a proof of Theorem \ref{thm.perm.len.redword1}.

\begin{proof}
[Proof of Theorem \ref{thm.perm.len.redword1} (sketched).]\textbf{(a)} (See
\cite[Exercise 5.2 \textbf{(e)}]{detnotes} for details.)

We proceed by induction on $\ell\left(  \sigma\right)  $:

\textit{Induction base:} If $\ell\left(  \sigma\right)  =0$, then
$\sigma=\operatorname*{id}$, so that we can write $\sigma$ as a composition of
$0$ simples.

\textit{Induction step:} Fix $h\in\mathbb{N}$. Assume (as the
IH\footnote{\textquotedblleft IH\textquotedblright\ means \textquotedblleft
induction hypothesis\textquotedblright.}) that Theorem
\ref{thm.perm.len.redword1} \textbf{(a)} holds for $\ell\left(  \sigma\right)
=h$.

Now, let $\sigma\in S_{n}$ be such that $\ell\left(  \sigma\right)  =h+1$. We
must prove that we can write $\sigma$ as a composition of $\ell\left(
\sigma\right)  $ simples.

We have $\ell\left(  \sigma\right)  =h+1>h\geq0$; hence, $\sigma$ has at least
one inversion. Thus, we cannot have $\sigma\left(  1\right)  \leq\sigma\left(
2\right)  \leq\cdots\leq\sigma\left(  n\right)  $. In other words, there
exists some $k\in\left[  n-1\right]  $ such that $\sigma\left(  k\right)
>\sigma\left(  k+1\right)  $. Let us fix such a $k$.

Lemma \ref{lem.perm.len.ssl} \textbf{(a)} yields%
\begin{align*}
\ell\left(  \sigma s_{k}\right)   &  =%
\begin{cases}
\ell\left(  \sigma\right)  +1, & \text{if }\sigma\left(  k\right)
<\sigma\left(  k+1\right)  ;\\
\ell\left(  \sigma\right)  -1, & \text{if }\sigma\left(  k\right)
>\sigma\left(  k+1\right)
\end{cases}
\\
&  =\ell\left(  \sigma\right)  -1\ \ \ \ \ \ \ \ \ \ \left(  \text{since
}\sigma\left(  k\right)  >\sigma\left(  k+1\right)  \right) \\
&  =h\ \ \ \ \ \ \ \ \ \ \left(  \text{since }\ell\left(  \sigma\right)
=h+1\right)  .
\end{align*}
Thus, the IH tells us that we can apply Theorem \ref{thm.perm.len.redword1}
\textbf{(a)} to $\sigma s_{k}$ instead of $\sigma$. We conclude that we can
write $\sigma s_{k}$ as a composition of $\ell\left(  \sigma s_{k}\right)  $
simples. In other words, we can write $\sigma s_{k}$ as a composition of $h$
simples (since $\ell\left(  \sigma s_{k}\right)  =h$). That is, we have%
\[
\sigma s_{k}=s_{i_{1}}s_{i_{2}}\cdots s_{i_{h}}\ \ \ \ \ \ \ \ \ \ \text{for
some }i_{1},i_{2},\ldots,i_{h}\in\left[  n-1\right]  .
\]
Consider these $i_{1},i_{2},\ldots,i_{h}$. Now,%
\[
\sigma=\underbrace{\sigma s_{k}}_{=s_{i_{1}}s_{i_{2}}\cdots s_{i_{h}}%
}\underbrace{s_{k}^{-1}}_{=s_{k}}=s_{i_{1}}s_{i_{2}}\cdots s_{i_{h}}s_{k}.
\]
This shows that we can write $\sigma$ as a composition of $h+1$ simples. In
other words, we can write $\sigma$ as a composition of $\ell\left(
\sigma\right)  $ simples (since $\ell\left(  \sigma\right)  =h+1$). Thus,
Theorem \ref{thm.perm.len.redword1} \textbf{(a)} holds for $\ell\left(
\sigma\right)  =h+1$. This completes the induction step, and so Theorem
\ref{thm.perm.len.redword1} \textbf{(a)} is proved by induction. \medskip

\textbf{(b)} (See \cite[Exercise 5.2 \textbf{(g)}]{detnotes} for details.)

We already know from Theorem \ref{thm.perm.len.redword1} \textbf{(a)} that we
can write $\sigma$ as a composition of $\ell\left(  \sigma\right)  $ simples.
It thus remains to show that we cannot write $\sigma$ as a composition of
fewer than $\ell\left(  \sigma\right)  $ simples.

This will clearly follow if we can show that%
\begin{equation}
\ell\left(  s_{i_{1}}s_{i_{2}}\cdots s_{i_{g}}\right)  \leq g
\label{pf.thm.perm.len.redword1.b.1}%
\end{equation}
for any $i_{1},i_{2},\ldots,i_{g}\in\left[  n-1\right]  $.

In order to prove (\ref{thm.perm.len.redword1}), we first make a simple
observation: For any $\sigma\in S_{n}$ and any $k\in\left[  n-1\right]  $, we
have%
\begin{equation}
\ell\left(  \sigma s_{k}\right)  \leq\ell\left(  \sigma\right)  +1.
\label{pf.thm.perm.len.redword1.b.2}%
\end{equation}
(This follows from Lemma \ref{lem.perm.len.ssl} \textbf{(a)}, since both
numbers $\ell\left(  \sigma\right)  +1$ and $\ell\left(  \sigma\right)  -1$
are $\leq\ell\left(  \sigma\right)  +1$.) Now, for any $i_{1},i_{2}%
,\ldots,i_{g}\in\left[  n-1\right]  $, we have%
\begin{align*}
&  \ell\left(  s_{i_{1}}s_{i_{2}}\cdots s_{i_{g}}\right) \\
&  \leq\ell\left(  s_{i_{1}}s_{i_{2}}\cdots s_{i_{g-1}}\right)
+1\ \ \ \ \ \ \ \ \ \ \left(  \text{by (\ref{pf.thm.perm.len.redword1.b.2}%
)}\right) \\
&  \leq\ell\left(  s_{i_{1}}s_{i_{2}}\cdots s_{i_{g-2}}\right)  +2\\
&  \ \ \ \ \ \ \ \ \ \ \left(  \text{since (\ref{pf.thm.perm.len.redword1.b.2}%
) yields }\ell\left(  s_{i_{1}}s_{i_{2}}\cdots s_{i_{g-1}}\right)  \leq
\ell\left(  s_{i_{1}}s_{i_{2}}\cdots s_{i_{g-2}}\right)  +1\right) \\
&  \leq\ell\left(  s_{i_{1}}s_{i_{2}}\cdots s_{i_{g-3}}\right)  +3\\
&  \ \ \ \ \ \ \ \ \ \ \left(  \text{since (\ref{pf.thm.perm.len.redword1.b.2}%
) yields }\ell\left(  s_{i_{1}}s_{i_{2}}\cdots s_{i_{g-2}}\right)  \leq
\ell\left(  s_{i_{1}}s_{i_{2}}\cdots s_{i_{g-3}}\right)  +1\right) \\
&  \leq\cdots\\
&  \leq\underbrace{\ell\left(  \operatorname*{id}\right)  }_{=0}%
+g\ \ \ \ \ \ \ \ \ \ \left(  \text{since the product of }0\text{ simples is
}\operatorname*{id}\right) \\
&  =g.
\end{align*}
This proves (\ref{pf.thm.perm.len.redword1.b.1}), and thus concludes our proof
of Theorem \ref{thm.perm.len.redword1} \textbf{(b)}.
\end{proof}

\begin{corollary}
\label{cor.perm.red.sigtau}Let $n\in\mathbb{N}$. \medskip

\textbf{(a)} We have $\ell\left(  \sigma\tau\right)  \equiv\ell\left(
\sigma\right)  +\ell\left(  \tau\right)  \operatorname{mod}2$ for all
$\sigma\in S_{n}$ and $\tau\in S_{n}$. \medskip

\textbf{(b)} We have $\ell\left(  \sigma\tau\right)  \leq\ell\left(
\sigma\right)  +\ell\left(  \tau\right)  $ for all $\sigma\in S_{n}$ and
$\tau\in S_{n}$. \medskip

\textbf{(c)} Let $k_{1},k_{2},\ldots,k_{q}\in\left[  n-1\right]  $, and let
$\sigma=s_{k_{1}}s_{k_{2}}\cdots s_{k_{q}}$. Then, $q\geq\ell\left(
\sigma\right)  $ and $q\equiv\ell\left(  \sigma\right)  \operatorname{mod}2$.
\end{corollary}

\begin{example}
Let $n=4$. Consider the two permutations $\sigma=3214$ and $\tau=3142$ (both
written in one-line notation). Then, $\ell\left(  \sigma\right)  =3$ and
$\ell\left(  \tau\right)  =3$. Now, the permutation $\sigma\tau=1342$ has
length $\ell\left(  \sigma\tau\right)  =2$.

Corollary \ref{cor.perm.red.sigtau} \textbf{(a)} says $\ell\left(  \sigma
\tau\right)  \equiv\ell\left(  \sigma\right)  +\ell\left(  \tau\right)
\operatorname{mod}2$. In other words, $2\equiv3+3\operatorname{mod}2$.

Corollary \ref{cor.perm.red.sigtau} \textbf{(b)} says $\ell\left(  \sigma
\tau\right)  \leq\ell\left(  \sigma\right)  +\ell\left(  \tau\right)  $. In
other words, $2\leq3+3$.
\end{example}

\begin{proof}
[Proof of Corollary \ref{cor.perm.red.sigtau} (sketched).]\textbf{(a)} (See
\cite[Exercise 5.2 \textbf{(b)}]{detnotes} for details.)

For any $\sigma\in S_{n}$ and any $k\in\left[  n-1\right]  $, we have%
\begin{equation}
\ell\left(  \sigma s_{k}\right)  \equiv\ell\left(  \sigma\right)
+1\operatorname{mod}2. \label{pf.cor.perm.red.sigtau.a.1}%
\end{equation}
(This follows from Lemma \ref{lem.perm.len.ssl} \textbf{(a)}, since both
numbers $\ell\left(  \sigma\right)  +1$ and $\ell\left(  \sigma\right)  -1$
are congruent to $\ell\left(  \sigma\right)  +1$ modulo $2$.)

Now, let $\sigma\in S_{n}$ and $\tau\in S_{n}$. Theorem
\ref{thm.perm.len.redword1} \textbf{(a)} yields that we can write $\tau$ as a
composition of $\ell\left(  \tau\right)  $ simples. In other words, we can
write $\tau$ as $\tau=s_{k_{1}}s_{k_{2}}\cdots s_{k_{q}}$ for some
$k_{1},k_{2},\ldots,k_{q}\in\left[  n-1\right]  $, where $q=\ell\left(
\tau\right)  $. Consider these $k_{1},k_{2},\ldots,k_{q}$. Now,%
\begin{align*}
\ell\left(  \sigma\tau\right)   &  =\ell\left(  \sigma s_{k_{1}}s_{k_{2}%
}\cdots s_{k_{q}}\right)  \ \ \ \ \ \ \ \ \ \ \left(  \text{since }%
\tau=s_{k_{1}}s_{k_{2}}\cdots s_{k_{q}}\right) \\
&  \equiv\ell\left(  \sigma s_{k_{1}}s_{k_{2}}\cdots s_{k_{q-1}}\right)
+1\ \ \ \ \ \ \ \ \ \ \left(  \text{by (\ref{pf.cor.perm.red.sigtau.a.1}%
)}\right) \\
&  \equiv\ell\left(  \sigma s_{k_{1}}s_{k_{2}}\cdots s_{k_{q-2}}\right)  +2\\
&  \ \ \ \ \ \ \ \ \ \ \left(  \text{since (\ref{pf.cor.perm.red.sigtau.a.1})
yields }\ell\left(  \sigma s_{k_{1}}s_{k_{2}}\cdots s_{k_{q-1}}\right)
\equiv\ell\left(  \sigma s_{k_{1}}s_{k_{2}}\cdots s_{k_{q-2}}\right)
+1\operatorname{mod}2\right) \\
&  \equiv\ell\left(  \sigma s_{k_{1}}s_{k_{2}}\cdots s_{k_{q-3}}\right)  +3\\
&  \ \ \ \ \ \ \ \ \ \ \left(  \text{since (\ref{pf.cor.perm.red.sigtau.a.1})
yields }\ell\left(  \sigma s_{k_{1}}s_{k_{2}}\cdots s_{k_{q-2}}\right)
\equiv\ell\left(  \sigma s_{k_{1}}s_{k_{2}}\cdots s_{k_{q-3}}\right)
+1\operatorname{mod}2\right) \\
&  \equiv\cdots\\
&  \equiv\ell\left(  \sigma\right)  +\underbrace{q}_{=\ell\left(  \tau\right)
}=\ell\left(  \sigma\right)  +\ell\left(  \tau\right)  \operatorname{mod}2.
\end{align*}
This proves Corollary \ref{cor.perm.red.sigtau} \textbf{(a)}. \medskip

\textbf{(b)} (See \cite[Exercise 5.2 \textbf{(c)}]{detnotes} for details.)

This is analogous to the proof of Corollary \ref{cor.perm.red.sigtau}
\textbf{(a)} (but using inequalities instead of congruences, and using
(\ref{pf.thm.perm.len.redword1.b.2}) instead of
(\ref{pf.cor.perm.red.sigtau.a.1})). \medskip

\textbf{(c)} Let $k_{1},k_{2},\ldots,k_{q}\in\left[  n-1\right]  $, and let
$\sigma=s_{k_{1}}s_{k_{2}}\cdots s_{k_{q}}$. We must prove that $q\geq
\ell\left(  \sigma\right)  $ and $q\equiv\ell\left(  \sigma\right)
\operatorname{mod}2$.

From $\sigma=s_{k_{1}}s_{k_{2}}\cdots s_{k_{q}}$, we obtain%
\begin{align*}
\ell\left(  \sigma\right)   &  =\ell\left(  s_{k_{1}}s_{k_{2}}\cdots s_{k_{q}%
}\right) \\
&  \equiv\ell\left(  s_{k_{1}}s_{k_{2}}\cdots s_{k_{q-1}}\right)
+1\ \ \ \ \ \ \ \ \ \ \left(  \text{by (\ref{pf.cor.perm.red.sigtau.a.1}%
)}\right) \\
&  \equiv\ell\left(  s_{k_{1}}s_{k_{2}}\cdots s_{k_{q-2}}\right)  +2\\
&  \ \ \ \ \ \ \ \ \ \ \left(  \text{since (\ref{pf.cor.perm.red.sigtau.a.1})
yields }\ell\left(  s_{k_{1}}s_{k_{2}}\cdots s_{k_{q-1}}\right)  \equiv
\ell\left(  s_{k_{1}}s_{k_{2}}\cdots s_{k_{q-2}}\right)  +1\operatorname{mod}%
2\right) \\
&  \equiv\ell\left(  s_{k_{1}}s_{k_{2}}\cdots s_{k_{q-3}}\right)  +3\\
&  \ \ \ \ \ \ \ \ \ \ \left(  \text{since (\ref{pf.cor.perm.red.sigtau.a.1})
yields }\ell\left(  s_{k_{1}}s_{k_{2}}\cdots s_{k_{q-2}}\right)  \equiv
\ell\left(  s_{k_{1}}s_{k_{2}}\cdots s_{k_{q-3}}\right)  +1\operatorname{mod}%
2\right) \\
&  \equiv\cdots\\
&  \equiv\underbrace{\ell\left(  \operatorname*{id}\right)  }_{=0}%
+\,q\ \ \ \ \ \ \ \ \ \ \left(  \text{since the product of }0\text{ simples is
}\operatorname*{id}\right) \\
&  =q\operatorname{mod}2.
\end{align*}
In other words, $q\equiv\ell\left(  \sigma\right)  \operatorname{mod}2$. A
similar argument (but using inequalities instead of congruences, and using
(\ref{pf.thm.perm.len.redword1.b.2}) instead of
(\ref{pf.cor.perm.red.sigtau.a.1})) shows that $q\geq\ell\left(
\sigma\right)  $. Thus, Corollary \ref{cor.perm.red.sigtau} \textbf{(c)} is proved.
\end{proof}

\begin{corollary}
\label{cor.perm.generated}Let $n\in\mathbb{N}$. Then, the group $S_{n}$ is
generated by the simples $s_{1},s_{2},\ldots,s_{n-1}$.
\end{corollary}

\begin{proof}
This follows directly from Theorem \ref{thm.perm.len.redword1} \textbf{(a)}.
\end{proof}

Theorem \ref{thm.perm.len.redword1} \textbf{(a)} shows that every permutation
$\sigma\in S_{n}$ can be represented as a product of $\ell\left(
\sigma\right)  $ simples (and in most cases, this can be done in many
different ways). It turns out that there is a rather explicit way to find such
a representation:

\begin{remark}
\label{rmk.perm.redword-lehmer}Let $n\in\mathbb{N}$ and $\sigma\in S_{n}$. Let
us represent the Lehmer code of $\sigma$ visually as follows:

We draw an (empty) $n\times n$-matrix.

For each $i\in\left[  n\right]  $, we put a cross $\times$ into the cell
$\left(  i,\sigma\left(  i\right)  \right)  $ of the matrix.

In the following, I will use the case $n=6$ and $\sigma=513462$ (in one-line
notation) as a running example. In this case, the matrix looks as follows:%
\[%
% [tikz figure removed]
%BeginExpansion
% [tikz figure removed]%
%EndExpansion
\]


Now, starting from each $\times$, we draw a vertical ray downwards and a
horizontal ray eastwards. I will call these two rays the \emph{Lehmer lasers}.
Here is how the rays look like in our running example:%
\[%
% [tikz figure removed]
%BeginExpansion
% [tikz figure removed]%
%EndExpansion
\]


This picture is called the \emph{Rothe diagram} of $\sigma$.

Explicitly, a cell $\left(  i,j\right)  $ has a $\circ$ in it if and only if%
\[
\sigma\left(  i\right)  >j\text{ and }\sigma^{-1}\left(  j\right)  >i
\]
(indeed, the vertical laser in column $j$ hits cell $\left(  i,j\right)  $ if
and only if $\sigma^{-1}\left(  j\right)  \leq i$, whereas the horizontal
laser in row $i$ hits cell $\left(  i,j\right)  $ if and only if
$\sigma\left(  i\right)  \leq j$).

If we substitute $\sigma\left(  j\right)  $ for $j$ in this statement, then we
obtain the following: A cell $\left(  i,\sigma\left(  j\right)  \right)  $ has
a $\circ$ in it if and only if%
\[
\sigma\left(  i\right)  >\sigma\left(  j\right)  \text{ and }j>i.
\]
In other words, a cell $\left(  i,\sigma\left(  j\right)  \right)  $ has a
$\circ$ in it if and only if $\left(  i,j\right)  $ is an inversion of
$\sigma$.

Thus,%
\[
\ell\left(  \sigma\right)  =\left(  \text{\# of inversions of }\sigma\right)
=\left(  \text{\# of }\circ\text{'s}\right)  ,
\]
and%
\[
\ell_{i}\left(  \sigma\right)  =\left(  \text{\# of }\circ\text{'s in row
}i\right)  \ \ \ \ \ \ \ \ \ \ \text{for each }i\in\left[  n\right]  .
\]


Finally, let us label the $\circ$'s as follows: For each $i\in\left[
n\right]  $, we label the $\circ$'s in row $i$ from right to left by the
simple transpositions $s_{i},s_{i+1},s_{i+2},\ldots,s_{i^{\prime}-1}$ where
$i^{\prime}=i+\ell_{i}\left(  \sigma\right)  $. (This works, since there are
precisely $\ell_{i}\left(  \sigma\right)  =i^{\prime}-i$ many $\circ$'s in row
$i$.) Here is how this labeling looks in our running example:%
\[%
% [tikz figure removed]
%BeginExpansion
% [tikz figure removed]%
%EndExpansion
\]


Finally, read the matrix row by row, starting with the top row, and reading
each row from left to right. The result, in our running example, is%
\[
s_{4}s_{3}s_{2}s_{1}s_{3}s_{4}s_{5}.
\]
Reading this as a product, we obtain a product of $\ell\left(  \sigma\right)
$ simples that equals $\sigma$.
\end{remark}

The claim of Remark \ref{rmk.perm.redword-lehmer} can be restated in a more
direct (if less visual) fashion:

\begin{proposition}
\label{prop.perm.redword-lehmer}Let $n\in\mathbb{N}$. Let $\sigma\in S_{n}$.
For each $i\in\left[  n\right]  $, we set%
\begin{equation}
a_{i}:=\operatorname*{cyc}\nolimits_{i^{\prime},i^{\prime}-1,i^{\prime
}-2,\ldots,i}=s_{i^{\prime}-1}s_{i^{\prime}-2}s_{i^{\prime}-3}\cdots s_{i},
\label{eq.prop.perm.redword-lehmer.1}%
\end{equation}
where $i^{\prime}=i+\ell_{i}\left(  \sigma\right)  $. Then, $\sigma=a_{1}%
a_{2}\cdots a_{n}$. (The second equality sign in
(\ref{eq.prop.perm.redword-lehmer.1}) is not hard to check. Note that
$a_{n}=\operatorname*{id}$.)
\end{proposition}

We refer to \cite[Exercise 5.21 parts \textbf{(b)} and \textbf{(c)}]{detnotes}
for a (detailed, but annoyingly long) proof of Proposition
\ref{prop.perm.redword-lehmer}. (You are probably better off proving it yourself.)
